\section{Methodology}
\label{sec:methodology}

\subsection{System overview}

\begin{figure*}[t]
  \centering
  \includegraphics[width=\textwidth]{../vowel_framework/vowel_framework.pdf}
  \caption{U-VOWEL execution loop. The VLA supplies the default action chunk
  and a branch-specific failure score. A VLA alarm makes the predictive branch
  compute a fresh video-conditioned context and decode an alternative action
  chunk from the current state. The visual frames represent the V2W training
  target; the deployed branch passes an intermediate representation to W2A
  without decoding a pixel-space rollout. The two scores are evaluated against
  separate thresholds and are never compared numerically.}
  \label{fig:vowel-framework}
\end{figure*}

U-VOWEL connects two independently trained action-generating policies through
the interface in Eq.~\eqref{eq:policy-query}. Both branches receive the same
language instruction from a synchronized decision state, but retain their
native observations, action representations and generative processes. The
direct VLA forms the default path. The video-world-model policy is queried only
when the VLA monitor raises an alarm. Figure~\ref{fig:vowel-framework} follows
one such decision from observation to execution.

Each policy first produces an action proposal and uncertainty measurements tied
to its own prediction process. A branch-specific temporal model converts those
measurements into a failure score. The controller then accepts one proposal,
shortens VLA execution or temporarily suppresses repeated calls to the
predictive branch. This separation prevents predicted variance and sample
disagreement from being treated as calibrated probabilities of task failure.

\subsection{Direct VLA policy}

\subsubsection{Action generation}

The direct branch modifies SimVLA~\cite{simvla}. A frozen SmolVLM-500M encoder
processes the language instruction and current agent and wrist RGB views. A
flow-matching action transformer combines this representation with an
eight-dimensional proprioceptive state and generates a normalized chunk of
\(H=10\) seven-dimensional Cartesian actions. For an expert chunk
\(A^\star\), Gaussian noise \(\epsilon\) and flow time
\(\gamma=0.001+0.999\tilde{\gamma}\), with
\(\tilde{\gamma}\sim\operatorname{Beta}(1.5,1)\), training constructs
\begin{equation}
  X_\gamma=(1-\gamma)A^\star+\gamma\epsilon,
  \qquad
  V^\star_\gamma=\epsilon-A^\star .
  \label{eq:vla-flow}
\end{equation}
The action transformer predicts the velocity field from
\((X_\gamma,\gamma)\), then inference integrates that field from noise to an
action chunk. The vision-language encoder remains frozen during uncertainty
fine-tuning; the action transformer and its output heads are trainable.

\subsubsection{Heteroscedastic velocity uncertainty}

We replace the original velocity output with parallel velocity and variance
heads. For action index \(i\), the second head emits \(r_{\theta,i}\), which is
mapped to a positive conditional variance,
\begin{equation}
  \sigma^2_{\theta,i}
  =\operatorname{softplus}(r_{\theta,i})+\varepsilon_{\mathrm{n}},
  \qquad \varepsilon_{\mathrm{n}}=10^{-6}.
  \label{eq:vla-variance}
\end{equation}
The modified policy is fine-tuned with an element-wise Beta-NLL objective
\cite{betanll}:
\begin{equation}
  \begin{aligned}
    \mathcal{L}_{\vla}
    &=
    \frac{1}{|\mathcal I|}
    \sum_{i\in\mathcal I}
    \operatorname{sg}\!\left[
      \left(\sigma^2_{\theta,i}\right)^\beta
    \right] \\
    &\quad{}\times
    \left[
      \frac{(V_{\theta,i}-V^\star_{\gamma,i})^2}
           {2\sigma^2_{\theta,i}}
      +\frac{1}{2}\log\sigma^2_{\theta,i}
      +\kappa
    \right],
    \qquad \beta=0.5 .
  \end{aligned}
  \label{eq:vla-beta-nll}
\end{equation}
Here \(\operatorname{sg}\) stops gradients through the Beta-NLL weight,
\(\mathcal I\) spans action timesteps and dimensions and
\(\kappa=-\tfrac{1}{2}\log\varepsilon_{\mathrm{n}}\). This constant keeps the
bracketed term nonnegative for
\(\sigma^2_{\theta,i}\geq\varepsilon_{\mathrm{n}}\). The residual term weights
velocity error by predicted precision, while the logarithmic term penalizes
indiscriminate variance inflation. The detached Beta-NLL factor reduces the
attenuation of high-variance examples without adding another gradient path
through the variance weight.

At inference, variance is retained along the Euler trajectory. A diagonal
first-order approximation accumulates
\begin{equation}
  P_{\gamma+\Delta\gamma}
  =P_\gamma+(\Delta\gamma)^2\sigma_\theta^2(X_\gamma,\gamma).
  \label{eq:vla-path-variance}
\end{equation}
Each query draws nine independently seeded \(H=10\) chunks from one encoded
observation. Candidate zero is the main proposal; eight alternatives provide
uncertainty evidence. A 49-dimensional record combines the main candidate's
last-step and path-integrated variance, denoising dynamics, chunk geometry and
dispersion across all nine samples. Separately, a seven-dimensional candidate-
disagreement descriptor is computed from the eight normalized alternatives:
centered and pairwise chunk distances, overall dispersion and translation,
rotation, gripper and flattened-coordinate dispersion. This descriptor is not
FIPER's entropy-based ACE. The deployed monitor retains eight of the 49
dimensions through a ranking performed on seen train/validation data, covering
initial denoising magnitude, update-direction flips, path variance and action-
sample dispersion. None of these inputs contains a task identifier, reward or
future observation.

\subsection{Predictive video--action policy}

\subsubsection{V2W--W2A generation}

The predictive branch separates a video-trained representation from action
decoding. Its video-to-world component (V2W) adapts a Cosmos Predict2 backbone
through rank-256 LoRA and trainable adaptive-normalization parameters
\cite{lora,cosmospolicy}. V2W is trained to predict future video from five
\(480\times640\) agent-view frames sampled from recent observation history and
the language embedding. The deployed policy does not complete or decode that
visual rollout. It stops at the first V2W denoising evaluation and extracts a
hidden representation after transformer block 20 as the predictive context
\(C_t\). The world-to-action component (W2A) is a separate diffusion
transformer trained while V2W remains frozen. It conditions on \(C_t\), the
ten-dimensional robot state and the video noise level, then generates 60
ten-dimensional tokens. Translation, six-dimensional rotation and gripper
outputs are decoded to the seven-dimensional environment action.

W2A uses the interpolation and velocity target in
Eq.~\eqref{eq:vla-flow}, with
\(\gamma=0.001+0.999\tilde{\gamma}\) and
\(\tilde{\gamma}\sim\operatorname{Beta}(1,1)\), and minimizes
\begin{equation}
  \mathcal{L}_{\mathrm{W2A}}
  =
  10\,\mathbb{E}_{\gamma,\epsilon}
  \left[
    \left\|
      V_\omega(X_\gamma,\gamma;q_t,C_t,\sigma)
      -V^\star_\gamma
    \right\|_2^2
  \right],
  \label{eq:w2a-loss}
\end{equation}
where \(\sigma\) is the V2W noise level. During training, the observation
condition is dropped with probability \(0.2\). This follows the flow-matching
formulation~\cite{flowmatching}. Although W2A predicts a long sequence, the
controller executes at most ten actions before observing the environment
again.

Repeated W2A samples expose action uncertainty without repeating the expensive
video pass. In the K1 configuration, one V2W context produces eight W2A
candidates. Their per-step and per-dimension variance, pairwise mean-squared
distance, receding-plan disagreement and denoising dynamics form the K1
uncertainty vector. Candidate zero remains the executable proposal; the other
samples provide evidence and do not compete for control. K3 instead generates
three V2W contexts. The main context produces eight W2A candidates and each
additional context produces one fixed-seed candidate, giving ten action chunks
for measuring sensitivity to the sampled V2W representation. Thus K1 and K3
name the number of V2W contexts, not the number of executed chunks or the
monitor history.

\subsubsection{Residual uncertainty in V2W}

Separately, we train a lightweight residual-variance head on the frozen V2W
backbone. Let \(h_{20,i}\) be a future token at block 20 and let
\(\rho_i^2\) be the mean squared clean-latent residual over its 16 channels. A
calibration table supplies a positive base variance
\(b_{q(\sigma)}\) for the current video-noise bin. The head predicts a
token-wise correction from \(h_{20,i}\) and an embedding of \(\log\sigma\):
\begin{equation}
  \begin{aligned}
    \ell_i &=
    \operatorname{clip}\!\left(
      \log b_{q(\sigma)}
      +\Delta_\phi(h_{20,i},\log\sigma),-8,5
    \right),\\
    v_i&=\exp(\ell_i).
  \end{aligned}
  \label{eq:v2w-variance}
\end{equation}
Only the variance head and noise embedding are optimized. With \(w_i=1\) for
future tokens and zero for observed conditioning tokens, the loss is
\begin{equation}
  \mathcal{L}_{\mathrm{V2W\mbox{-}NLL}}
  =
  \frac{1}{2}
  \frac{\sum_i w_i\left(\rho_i^2/v_i+\ell_i\right)}
       {\sum_i w_i+\varepsilon_{\mathrm{n}}}.
  \label{eq:v2w-nll}
\end{equation}
This head estimates conditional latent-residual scale. It does not by itself
measure task failure or complete epistemic uncertainty. Its summaries enter
the uncertainty-feature ablations. The selected K1 controller uses the
action-uncertainty profile described above and therefore does not depend on
V2W variance-head features.

\subsection{Branch-specific temporal risk}

Let \(R_t^b\) denote the normalized query record built for branch \(b\). The
record may contain sequential action or uncertainty tokens, a bounded history
and current summary features. Each monitor computes
\begin{equation}
  s_t^b=\operatorname{sigmoid}\!\left(f_b(R_t^b)\right),
  \qquad b\in\{\vla,\wm\}.
  \label{eq:risk-score}
\end{equation}
The records and functions are branch-specific because the two generative
models expose different evidence.

\subsubsection{VLA temporal risk monitor}

The VLA monitor takes FIPER~\cite{fiper} as a methodological starting point.
FIPER accumulates RND-OE observation-embedding novelty and grid-based
action-chunk entropy (ACE) over short windows, calibrates both signals from
successful rollouts and combines their alarms by a logical AND. Our VLA monitor
instead learns one supervised score from seen successful and failed rollouts,
using candidate-action evidence and its recent history. Its seven-value
candidate-disagreement descriptor is not FIPER ACE, and it uses neither RND nor
direct observation embeddings.

The VLA monitor uses a late-fusion Transformer. Inputs are standardized using
statistics fitted only on the training split. Its sequence path contains 16
history tokens of dimension 21---proprioception (8), the previous executed
action (7) and six candidate-disagreement values---followed by ten 7D tokens
from the normalized main proposal. Separate linear projections map both token
types to width 128. A learned CLS token and positional embeddings precede a
three-layer, four-head Transformer with 512-dimensional feed-forward blocks.
Startup histories are zero-padded and processed without an attention mask.

The static 51D path concatenates 28 main-chunk statistics (first, mean,
standard deviation and endpoint change), the seven candidate-disagreement
values, current 8D proprioception and the selected eight uncertainty features.
A linear--GELU projection maps it to width 128. Its output is concatenated with
the Transformer CLS output, normalized and passed through a 128D hidden layer
to the risk logit. Neither task identity nor timestep is an input.

The VLA monitor is trained as a binary classifier from rollout outcomes. A query
in a failed rollout receives a failure target even when that query did not
cause the eventual failure. Let \(F_e=1-Y_e\) denote episode failure, and let
\(N_0\) and \(N_1\) be the numbers of successful and failed training rows. Its
positive-class weight is
\(p=\max(1,N_0/\max(1,N_1))\), and the implemented objective is
\begin{equation}
  \mathcal{L}_{\mathrm{risk}}^{\vla}
  =-\frac{1}{N}\sum_{e,t}
  \left[
    pF_e\log s_{e,t}^{\vla}
    +(1-F_e)\log(1-s_{e,t}^{\vla})
  \right].
  \label{eq:risk-loss}
\end{equation}
Consequently, \(s_t^\vla\) is an empirical warning score for the observed
episode prefix, not a causal probability that executing \(A_t^\vla\) will fail.

VLA operating points are computed on a held-out seen validation split: one
maximizes row-level F1 using both classes, while success-only alternatives use
empirical quantiles of successful row scores. The operating point is selected
according to the desired seen false-alarm/detection trade-off and frozen before
unseen evaluation. An episode alarms if any query crosses it. This is empirical
row-score calibration, not an episode-level conformal guarantee.

\subsubsection{World-model temporal risk monitor}

The selected K1 world-model monitor uses an eight-query history and a
two-layer GRU. At each query it receives 74 scalar W2A uncertainty features,
including candidate variance, pairwise disagreement, plan overlap and
denoising summaries. Sixteen temporal action-uncertainty tokens preserve the
per-step structure of those statistics. The GRU and static branch have width
128 and their fused representation passes through a 64-dimensional latent
layer to the risk logit. V2W residuals, multi-context disagreement and
candidate centrality define richer profiles used in the ablations, but are
absent from this K1 feature contract.

The world-model threshold is calibrated from the episode maximum score on
successful held-out rollouts. For \(n\) calibration episodes with maxima
\(g_e=\max_t s_{e,t}^{\wm}\), its level-\(\alpha\) threshold is the corrected
empirical quantile
\begin{equation}
  \tau_\wm(\alpha)
  =
  g_{(\min\{n,\lceil(n+1)(1-\alpha)\rceil\})},
  \label{eq:wm-conformal-threshold}
\end{equation}
where \(g_{(k)}\) is the \(k\)-th order statistic. This controls the marginal
episode-level false-alarm rate under the usual exchangeability assumption
\cite{conformal}; it does not guarantee calibration after arbitrary sequential
distribution shift.

\subsection{Selective execution}

Algorithm~\ref{alg:uvowel} combines the two monitors without comparing their
score magnitudes. The state variable \(c\) is the number of environment actions
remaining before the world-model branch becomes eligible after a rejected
proposal. When \(c=0\), a VLA alarm triggers a fresh V2W--W2A query from the
same evidence \(E_t\). An accepted world-model proposal supplies the next
action prefix. If its own monitor also alarms, the proposal is discarded, the
controller executes a shorter prefix of the already available VLA chunk and
starts the latch. The selected K1 controller uses \(H=10\), \(h=5\) and
\(C=50\).

\begin{algorithm}[t]
  \caption{\method{} selective execution for one episode}
  \label{alg:uvowel}
  \footnotesize
  \begin{algorithmic}[1]
    \Require policies \(\pi_\vla,\pi_\wm\), monitors \(f_\vla,f_\wm\),
      thresholds \(\tau_\vla,\tau_\wm\), horizons \(H>h\), latch span \(C\)
    \State \(c\gets0\) \Comment{remaining ineligible environment actions}
    \While{episode active and action budget remains}
      \State acquire current evidence \(E_t\)
      \State \((A^\vla,u^\vla)\gets\pi_\vla(E_t)\);
        \(s^\vla\gets f_\vla(R_t^\vla)\)
      \If{\(s^\vla\leq\tau_\vla\)}
        \State \(m\gets H\) if \(c=0\), else \(\min(H,c)\)
        \State \(n\gets\Call{ExecutePrefix}{A^\vla,m}\);
          \(c\gets\max(c-n,0)\)
      \ElsIf{\(c>0\)}
        \State \(n\gets\Call{ExecutePrefix}{A^\vla,\min(h,c)}\);
          \(c\gets\max(c-n,0)\)
      \Else
        \State \((A^\wm,u^\wm)\gets\pi_\wm(E_t)\)
          \Comment{fresh query from the current state}
        \State \(s^\wm\gets f_\wm(R_t^\wm)\)
        \If{\(s^\wm\leq\tau_\wm\)}
          \State \(\Call{ExecutePrefix}{A^\wm,H}\)
        \Else
          \State \(\Call{DiscardProposal}{A^\wm}\)
          \State \(c\gets C\)
          \State \(n\gets\Call{ExecutePrefix}{A^\vla,\min(h,c)}\);
            \(c\gets\max(c-n,0)\)
        \EndIf
      \EndIf
    \EndWhile
  \end{algorithmic}
\end{algorithm}

While the latch is active, the VLA is still queried and scored at every
decision point. A high VLA score no longer invokes V2W--W2A and receives the
short horizon \(h\). A low score may use the ordinary horizon \(H\), capped at
latch expiry so that eligibility is reconsidered at the next state. The latch
is decremented by actions actually applied, including early termination within
a chunk. Once it expires, the next VLA alarm runs new world-model inference.
The rejected action buffer is never reused; its score remains only in the
monitor history and audit log.

\FloatBarrier
